% ch37_simplex_category.tex — Volume IV Chapter 1
% Retrofitted with full Vol I-III pedagogy: cycle stages, etymology, dialogs.

\chapter{The Simplex Category $\Delta$}

\begin{overviewbox}
Volume~IV opens on the combinatorial substrate that every model of
$\infty$-category theory rests upon. Before we can speak of quasi-categories
(Chapter~43), simplicial categories, or complete Segal spaces, we need the
indexing category $\Delta$ — the category of finite non-empty totally
ordered sets and order-preserving maps between them.

We approach $\Delta$ from three angles, walking the cycle at a gentle
pace. \emph{Geometrically}, each object $[n] = \{0 < 1 < \cdots < n\}$ is
the vertex set of an $n$-dimensional simplex (point, line segment,
triangle, tetrahedron, \ldots), and the morphisms are the combinatorial
shadows of affine maps between simplices that preserve vertex orderings.
\emph{Combinatorially}, the morphisms are generated by elementary
\emph{face maps} $\delta^i$ (one-point omissions) and \emph{degeneracy
maps} $\sigma^i$ (one-point duplications) subject to a small list of
relations, the cosimplicial identities. \emph{Universally}, the augmented
simplex category $\Delta_+$ is the free strict monoidal category on a
single monoid object — a remarkable theorem of André Joyal that explains
why $\Delta$ appears \emph{everywhere a monoid acts}.

The next two chapters use $\Delta$ to build first $\mathbf{sSet} =
\mathbf{Set}^{\Delta^{\mathrm{op}}}$ (the category of \emph{simplicial sets},
Chapter~38) and then the geometric realization $\lvert{-}\rvert :
\mathbf{sSet} \to \mathbf{Top}$ (Chapter~39). The chapter after that
turns to model-categorical machinery (Chapters~40--42). The journey to
quasi-categories begins here, with monotone maps between ordered sets.
\end{overviewbox}


% ═══════════════════════════════════════════════════════════════════════════
\section{The Category $\Delta$}
\label{sec:delta-category}
% ═══════════════════════════════════════════════════════════════════════════

\subsection{Informally}

Before any definitions, fix a mental picture. Imagine a sequence of
objects, one for each dimension:
\begin{itemize}
  \item A \emph{point} (dimension 0): one vertex.
  \item A \emph{line segment} (dimension 1): two vertices, joined by an edge.
  \item A \emph{triangle} (dimension 2): three vertices, three edges,
        one $2$-dimensional interior.
  \item A \emph{tetrahedron} (dimension 3): four vertices, six edges,
        four triangular faces, one $3$-dimensional interior.
  \item And so on: an \emph{$n$-simplex} has $n+1$ vertices and is the
        ``simplest possible'' shape of dimension $n$.
\end{itemize}

Each of these shapes has a \emph{combinatorial skeleton}: the vertex
set, totally ordered (we always agree on which vertex is ``first,''
which is ``second,'' and so on). For an $n$-simplex with vertices
labeled $0, 1, \ldots, n$, the skeleton is simply the ordered set
$\{0 < 1 < \cdots < n\}$.

Now imagine ways to move between simplices that respect both the shape
and the ordering. An edge $\{0 < 1\}$ has three natural maps into a
triangle $\{0 < 1 < 2\}$: send both endpoints into a single vertex
(``collapse the edge to a point''), or include them as one of the
three edges of the triangle. A triangle has many natural maps into a
tetrahedron: include as a face, or collapse two vertices, or send
everything to a single vertex. All of these — collapses, inclusions,
and combinations — are exactly the \emph{order-preserving} (or
\emph{monotone}) maps between the underlying ordered vertex sets.

The category $\Delta$ collects every such map. Its objects are the
combinatorial skeletons; its morphisms are the order-preserving maps;
composition is the obvious composition of functions. The category
captures, in pure combinatorial form, every way that simplices fit
together.

\begin{nameorigin}
The symbol \term{$\Delta$} is the Greek capital letter \emph{delta},
chosen because the standard 2-simplex — a triangle — has shape $\Delta$.
The convention dates to Eilenberg and Steenrod's \emph{Foundations of
Algebraic Topology} (1952), where the simplex category was first
isolated as the indexing object for singular homology. The word
\term{simplex} is from Latin \emph{simplex} (``simple, single-fold''),
since an $n$-simplex is the simplest convex shape of dimension $n$.
\end{nameorigin}

\begin{readernote}
The single most useful thing to remember about $\Delta$ is that its
morphisms are abundant: between any two objects $[m]$ and $[n]$ there
are \emph{many} monotone maps (we will count them in
Example~\ref{ex:counting-monotone}). This contrasts with categories
like $\mathrm{posets}$ where you might expect maps to be rare. The
abundance is what makes $\Delta$ useful as an \emph{indexing} category:
many maps means many ways to ``compare'' simplicial data of different
dimensions.
\end{readernote}

\subsection{Expanding}

Three lenses through which the same category appears.

\paragraph{The combinatorial lens.}
$\Delta$ is the category of finite non-empty totally ordered sets.
Choosing one representative of each isomorphism class (the standard
shape $\{0 < 1 < \cdots < n\}$) gives the objects $[0], [1], [2],
\ldots$. The morphisms are exactly the monotone functions.

\paragraph{The geometric lens.}
$\Delta$ is the category of \emph{combinatorial simplices}. The object
$[n]$ corresponds to the standard $n$-simplex
$\lvert\Delta^n\rvert \subset \mathbb{R}^{n+1}$ (the convex hull of the
standard basis vectors). Each monotone map $f : [m] \to [n]$ extends
linearly to an affine map of simplices $\lvert f\rvert : \lvert\Delta^m\rvert
\to \lvert\Delta^n\rvert$ that sends the $i$-th vertex of the source to
the $f(i)$-th vertex of the target.

\paragraph{The algebraic lens.}
$\Delta$ is generated by two families of elementary maps — the face maps
$\delta^i$ (inclusions of $(n-1)$-faces) and degeneracy maps $\sigma^i$
(collapses of adjacent vertex pairs) — subject to a finite set of
relations called the \emph{cosimplicial identities}
(\S\ref{sec:face-degeneracy}). This is the generators-and-relations
presentation that makes $\Delta$ usable in software and explicit
computations.

\begin{dialog}
\textbf{Reader:} ``Why ordered sets, and not just sets?''

\textbf{Author:} The orderings give the simplices a notion of
``orientation.'' An ordered triangle $\{0 < 1 < 2\}$ knows that
``vertex 0 comes before vertex 1, which comes before vertex 2'' — and
this matters as soon as we want to define face maps without ambiguity
(``the face opposite vertex 2'') or degeneracy maps
(``collapse vertex 1 onto vertex 2''). Unordered simplices would
require us to keep track of all possible relabellings; the ordering
fixes a canonical form.

\textbf{Reader:} ``So $\Delta$ is rigid in a way ordinary set theory isn't.''

\textbf{Author:} Exactly. $\Delta$ is rigid \emph{combinatorially}: it
has no non-identity automorphisms (Exercise: prove this). Every
isomorphism $[n] \to [n]$ is the identity. This is the structural
payoff of the ordering.
\end{dialog}

\begin{example}{Three monotone maps in pictures}
Consider three monotone maps $[1] \to [2]$:
\begin{itemize}
  \item $0 \mapsto 0$, $1 \mapsto 1$: includes the edge as ``front'' of the triangle.
  \item $0 \mapsto 0$, $1 \mapsto 2$: includes the edge as the long ``base.''
  \item $0 \mapsto 0$, $1 \mapsto 0$: collapses the edge to a single vertex.
\end{itemize}
The first two are injections; the third is the constant. Each is a
distinct map. (See Example~\ref{ex:counting-monotone} for the full count.)
\end{example}

\subsection{Formalizing}

We now collect the picture into a precise definition.

\begin{definition}[The simplex category $\Delta$]
\label{def:simplex-category}
The \term{simplex category} $\Delta$ has:
\begin{itemize}
  \item \textbf{Objects.} For each integer $n \geq 0$, an object $[n] =
        \{0, 1, \ldots, n\}$ regarded as a totally ordered set under the
        usual $<$.
  \item \textbf{Morphisms.} A morphism $f : [m] \to [n]$ is a function
        $f : \{0, \ldots, m\} \to \{0, \ldots, n\}$ that is \term{monotone}:
        $i \leq j$ in $[m]$ implies $f(i) \leq f(j)$ in $[n]$.
  \item \textbf{Composition.} Function composition; the composite of two
        monotone maps is monotone.
  \item \textbf{Identities.} The identity function $\mathrm{id}_{[n]}$.
\end{itemize}
\end{definition}

\begin{howtosay}
$\Delta$ \sayit{``the simplex category'' or ``capital delta''}; $[n]$
\sayit{``$n$-bracket'' or ``the standard $n$-simplex (combinatorial)''}.
A monotone map $f : [m] \to [n]$ may be read aloud as \sayit{``$f$ from
$[m]$ to $[n]$'' or ``a monotone map $m$-to-$n$.''} Saying ``$[n]$ has
$n+1$ elements'' is a useful catch-phrase to keep the off-by-one
convention front-of-mind.
\end{howtosay}

\begin{nameorigin}
\term{Monotone} comes from Greek \emph{mónos} (single) + \emph{tónos}
(tone, tension, stretch). A monotone function ``moves in a single
direction'' — never reversing. The mathematical usage dates to the
$19$th century in real analysis (\emph{monotonically increasing} sequences
of real numbers). In $\Delta$ we use it for functions of ordered finite
sets; the meaning is the same.
\end{nameorigin}

\begin{readernote}
The order-preservation condition $i \leq j \Rightarrow f(i) \leq f(j)$ is
\emph{weak} (it uses $\leq$, not $<$), so a monotone map may send two
distinct elements to the same image. This is what allows the
``collapse'' morphisms (the degeneracies) to exist in $\Delta$.
\end{readernote}

\begin{remark}[Cardinality vs.\ index]
Beware: $[n]$ has $n+1$ \emph{elements}, not $n$. The off-by-one is
geometric: an $n$-simplex (an $n$-dimensional shape) has $n+1$ vertices.
Authors who write $\boldsymbol{n} = \{0, 1, \ldots, n-1\}$ shift the
indexing by one but pay for it in formulas; we follow the geometric
convention throughout.
\end{remark}

\begin{remark}[Where $\Delta$ appears]
$\Delta$ is the indexing category for \emph{simplicial objects}:
contravariant functors $\Delta^{\mathrm{op}} \to \mathcal{C}$ for any
category $\mathcal{C}$. When $\mathcal{C} = \mathbf{Set}$ we get
\emph{simplicial sets} (Chapter~38); when $\mathcal{C} = \mathbf{Grp}$
we get \emph{simplicial groups} (classifying-space machinery); when
$\mathcal{C} = \mathbf{CommRing}$ we get \emph{simplicial commutative
rings} (one foundation of derived algebraic geometry, Chapter~58). In
each case, the combinatorial geometry lives in $\Delta$ and the values
live in the chosen target.
\end{remark}

\subsection{Formally: Small cases}

\begin{example_thm}[{The objects $[0]$ and $[1]$}]
$[0] = \{0\}$ is the one-element ordered set. It is a \emph{terminal}
object of $\Delta$: for every $[n]$ there is exactly one monotone map
$[n] \to [0]$ (the constant function $0$). Geometrically, $[0]$ is the
$0$-simplex — a single point.

$[1] = \{0 < 1\}$ has three monotone endomorphisms (the identity and the
two constants), and exactly two non-constant monotone maps $[0] \to [1]$
(picking out the element $0$ or the element $1$). $[1]$ is \emph{not}
initial in $\Delta$: there is not a unique map $[0] \to [1]$, but two —
and an initial object must have a \emph{unique} map to every other
object. Geometrically, $[1]$ is the $1$-simplex — a line segment.
\end{example_thm}

\begin{example_thm}[Counting monotone maps]
\label{ex:counting-monotone}
The number of monotone maps $[m] \to [n]$ equals $\binom{m + n + 1}{m + 1}$.
Equivalently: choose a non-decreasing sequence
\begin{equation*}
  0 \leq f(0) \leq f(1) \leq \cdots \leq f(m) \leq n.
\end{equation*}
This is the same as choosing a multiset of size $m+1$ from $\{0, 1, \ldots,
n\}$, equivalently $\binom{(n+1) + (m+1) - 1}{m+1} = \binom{m+n+1}{m+1}$
stars-and-bars combinations.
\end{example_thm}

\begin{example}{Spot-checking the count}
Monotone maps $[1] \to [2]$: $\binom{1 + 2 + 1}{1 + 1} = \binom{4}{2} = 6$.
Listing them: $(0,0), (0,1), (0,2), (1,1), (1,2), (2,2)$ — six pairs of
non-decreasing values. Of these, three are injections (rows $1, 2, 4$),
two are surjections-on-image (rows $1, 2, 4$ again — these inject), and
the remaining three are constants. The count checks out.
\end{example}

\begin{readernote}
$[0]$ is terminal but not initial in $\Delta$. The asymmetry is
characteristic: $\Delta$ has a terminal object (the one-point set) but
\emph{no initial object}, because the empty set is not present (it would
need cardinality zero, but $[n]$ has $n+1 \geq 1$ elements). The
augmented variant $\Delta_+$ in \S\ref{sec:universal} adjoins the empty
set $[-1] = \emptyset$ as an initial object — exactly what is needed to
make $\Delta_+$ strict-monoidal with $[-1]$ as unit.
\end{readernote}

\begin{remark}[$\Delta$ is small and locally finite]
$\Delta$ has countably many objects and, for each pair $[m], [n]$,
\emph{finitely} many morphisms (by Example~\ref{ex:counting-monotone}).
It is therefore a small category, with finite hom-sets — making $\Delta$
a useful indexing category for both set-valued diagrams and topological
diagrams.
\end{remark}


% ═══════════════════════════════════════════════════════════════════════════
\section{Face and Degeneracy Maps}
\label{sec:face-degeneracy}
% ═══════════════════════════════════════════════════════════════════════════

\subsection{Informally}

A monotone map $f : [m] \to [n]$ does one of three things:
\begin{itemize}
  \item Possibly \emph{omits} some elements of $[n]$ from its image.
  \item Possibly \emph{repeats} some images (sending two adjacent elements
        of $[m]$ to the same place in $[n]$).
  \item Otherwise: passes elements through in order.
\end{itemize}

It turns out that every monotone map factors uniquely as a sequence of
\emph{elementary omissions} (called face maps) followed by a sequence of
\emph{elementary repeats} (called degeneracy maps) — or vice versa,
depending on the factorization convention. The face and degeneracy maps
are the atomic building blocks of $\Delta$.

\begin{nameorigin}
\term{Face map} is geometric: in an $n$-simplex, the face opposite vertex
$i$ is the $(n-1)$-simplex obtained by deleting $i$ from the vertex set.
The face map $\delta^i : [n-1] \to [n]$ \emph{embeds} the $(n-1)$-simplex
into the $n$-simplex as that face.

\term{Degeneracy map} comes from Latin \emph{degenerare} (to depart from
one's lineage, to deteriorate). In mathematics ``degenerate'' generally
means ``collapsed to a simpler form'' (a degenerate triangle is one
where two vertices coincide). The degeneracy map $\sigma^i : [n+1] \to
[n]$ collapses two adjacent vertices, making the source $(n+1)$-simplex
``degenerate'' onto a non-degenerate face of the target $n$-simplex.
\end{nameorigin}

\subsection{Expanding}

The face map $\delta^i : [n-1] \to [n]$ ``inserts $i$ into $[n-1]$'s
image without using it.'' Concretely: it sends $0, 1, \ldots, i-1$ to
themselves, and $i, i+1, \ldots, n-1$ all shifted up by one to become
$i+1, i+2, \ldots, n$. The vertex $i \in [n]$ is the one omitted.

The degeneracy map $\sigma^i : [n+1] \to [n]$ ``duplicates the image
at $i$.'' Concretely: it sends $0, 1, \ldots, i$ to themselves, and
$i+1, i+2, \ldots, n+1$ all shifted down by one to become $i, i+1,
\ldots, n$. The fibre over $i \in [n]$ has two elements (namely $i$ and
$i+1$ in $[n+1]$); the fibres over every other element are singletons.

\begin{example}{Faces and degeneracies in low dimension}
\textbf{The maps $\delta^0, \delta^1 : [0] \to [1]$:}
$\delta^0$ skips $0 \in [1]$, so $0 \in [0]$ maps to $1 \in [1]$;
$\delta^1$ skips $1 \in [1]$, so $0 \in [0]$ maps to $0 \in [1]$.
Geometrically, $[0]$ is a point and $[1]$ is an edge: $\delta^0$ and
$\delta^1$ pick out the two endpoints.

\textbf{The map $\sigma^0 : [1] \to [0]$:}
The unique map $[1] \to [0]$ — both elements collapse onto the single
target. Geometrically, the edge collapses to a point.

\textbf{The maps $\delta^0, \delta^1, \delta^2 : [1] \to [2]$:}
The three edges of a triangle. $\delta^0$ skips $0$: gives the edge
$\{1, 2\}$ (opposite vertex $0$). $\delta^1$ skips $1$: gives the
``long'' edge $\{0, 2\}$. $\delta^2$ skips $2$: gives the edge $\{0, 1\}$.

\textbf{The maps $\sigma^0, \sigma^1 : [2] \to [1]$:}
Two ways to collapse a triangle onto an edge by identifying two adjacent
vertices. $\sigma^0$ identifies $\{0, 1\}$ in $[2]$ to $0 \in [1]$;
$\sigma^1$ identifies $\{1, 2\}$ in $[2]$ to $1 \in [1]$.
\end{example}

\begin{readernote}
\emph{Why the superscripts?} The convention $\delta^i$ and $\sigma^i$
(superscripts, not subscripts) emphasizes the \emph{covariant} direction:
the maps go $[n-1] \to [n]$ and $[n+1] \to [n]$, both in $\Delta$
itself. When we later pass to the opposite category $\Delta^{\mathrm{op}}$
to define simplicial sets, the induced face and degeneracy operators on
the simplicial set are written $d_i$ and $s_i$ (subscripts) — Chapter~38.
The convention helps signal which direction you are working in.
\end{readernote}

\subsection{Formalizing}

\begin{definition}[Face map]
\label{def:face-map}
For $0 \leq i \leq n$, the \term{$i$-th face map}
\begin{equation*}
  \delta^i : [n-1] \longrightarrow [n]
\end{equation*}
is the unique \emph{injective} monotone map whose image omits exactly the
element $i \in [n]$:
\begin{equation*}
  \delta^i(j) \;=\; \begin{cases} j & \text{if } j < i, \\ j + 1 & \text{if }
  j \geq i. \end{cases}
\end{equation*}
\end{definition}

\begin{definition}[Degeneracy map]
\label{def:degeneracy-map}
For $0 \leq i \leq n$, the \term{$i$-th degeneracy map}
\begin{equation*}
  \sigma^i : [n+1] \longrightarrow [n]
\end{equation*}
is the unique \emph{surjective} monotone map whose fibre over $i$ has two
elements (namely $i$ and $i+1$):
\begin{equation*}
  \sigma^i(j) \;=\; \begin{cases} j & \text{if } j \leq i, \\ j - 1 & \text{if }
  j > i. \end{cases}
\end{equation*}
\end{definition}

\begin{howtosay}
$\delta^i$ \sayit{``face delta-$i$''} ``skips $i$''; $\sigma^i$
\sayit{``degeneracy sigma-$i$''} ``doubles $i$.'' Some texts write $d_i$
and $s_i$ instead; we follow the more common ``cosimplicial'' (covariant)
convention with $\delta, \sigma$ throughout this chapter and reserve
$d_i, s_i$ for the contravariant simplicial-set operators of Chapter~38.
\end{howtosay}

\subsection{Reorganizing: the cosimplicial identities}

Face and degeneracy maps interact in a small, finite number of ways. The
package of relations is called the \term{cosimplicial identities}.

\begin{nameorigin}
\term{Cosimplicial} is the dual (``co-'') of \term{simplicial}: where a
\emph{simplicial} object is a contravariant functor on $\Delta$ (i.e., a
functor $\Delta^{\mathrm{op}} \to \mathcal{C}$), a \emph{cosimplicial}
object is a covariant functor $\Delta \to \mathcal{C}$. The face and
degeneracy maps live in $\Delta$ itself; their relations are the
\emph{cosimplicial} identities. The dual relations among $d_i, s_i$
(face and degeneracy \emph{operators} on a simplicial object) are
called the \term{simplicial identities} (Chapter~38).
\end{nameorigin}

\begin{theorem}[Cosimplicial identities]
\label{thm:cosimplicial-identities}
The face and degeneracy maps in $\Delta$ satisfy:
\begin{align*}
  \delta^j \delta^i &= \delta^i \delta^{j-1}      & & \text{if } i < j, \\
  \sigma^j \sigma^i &= \sigma^i \sigma^{j+1}      & & \text{if } i \leq j, \\
  \sigma^j \delta^i &= \delta^i \sigma^{j-1}      & & \text{if } i < j, \\
  \sigma^j \delta^i &= \mathrm{id}                 & & \text{if } i = j \text{ or } i = j+1, \\
  \sigma^j \delta^i &= \delta^{i-1} \sigma^j      & & \text{if } i > j+1.
\end{align*}
Every identity among composites of $\delta$'s and $\sigma$'s in $\Delta$
is a consequence of these five.
\end{theorem}

\begin{proof}[Proof sketch]
Each identity is checked by tracking what each side does to a generic
element. For example, the first: $\delta^j \delta^i : [n-2] \to [n]$
skips $i$ and then $j$ (now shifted); under the assumption $i < j$, the
``$j$ after $i$ has shifted'' becomes $j-1$ in the codomain $[n-1]$, and
the order $\delta^i \delta^{j-1}$ — skip $j-1$, then $i$ — produces the
same omission set $\{i, j\}$ with the same target indexing.

The completeness claim (every relation follows) is a consequence of the
epi-mono factorization in $\Delta$ (Section~\ref{sec:universal}).
\qed
\end{proof}

\begin{readernote}
The cosimplicial identities are best remembered as: \emph{faces commute
past faces with a shift; degeneracies commute past degeneracies with a
shift; faces and degeneracies either commute or annihilate}. The
annihilation case $\sigma^j \delta^i = \mathrm{id}$ for $i \in \{j, j+1\}$
is the key geometric content: \emph{doubling a vertex and then deleting
one copy of the doubled pair recovers the original vertex}.

A useful first practice: pick small $n$ and write out a few cases by
hand. The patterns are easy to see in dimension $n = 2$ or $3$, where
all maps are visualizable.
\end{readernote}

\subsection{Formally: generators}

\begin{theorem}[Generators of $\Delta$]
\label{thm:delta-generators}
Every morphism in $\Delta$ is a composite of face and degeneracy maps.
Specifically, any non-identity morphism $f : [m] \to [n]$ admits a unique
factorization
\begin{equation*}
  f \;=\; \delta^{i_1} \delta^{i_2} \cdots \delta^{i_p} \, \sigma^{j_1} \sigma^{j_2} \cdots \sigma^{j_q}
\end{equation*}
with $n \geq i_1 > i_2 > \cdots > i_p \geq 0$ and $0 \leq j_1 < j_2 < \cdots
< j_q < m$, where $p$ is the number of elements of $[n]$ not in the image
of $f$ and $q$ is the number of repetitions in the image (counting
multiplicity).
\end{theorem}

\begin{proof}[Proof sketch]
Every monotone map $f : [m] \to [n]$ factors uniquely as a surjection
followed by an injection (the standard epi-mono factorization, which
exists in $\Delta$ because monotone maps are determined by their images):
write $f = e \circ s$ where $s : [m] \twoheadrightarrow [k]$ is the
surjective part (collapsing repeated values) and $e : [k] \hookrightarrow
[n]$ is the injective part (the image inclusion).

Surjective monotone maps $[m] \twoheadrightarrow [k]$ are uniquely
composites of degeneracies in increasing index order; injective monotone
maps $[k] \hookrightarrow [n]$ are uniquely composites of faces in
decreasing index order. The unique factorization claim follows.
\qed
\end{proof}

\begin{remark}[Why the index order matters]
The ``decreasing $i$'s, increasing $j$'s'' convention is the unique
\emph{normal form}: any other ordering can be rewritten into this one
using the cosimplicial identities. The convention is one of the cleanest
in combinatorial topology — once you internalize it, you can read off
the structure of any monotone map at a glance.
\end{remark}

\begin{remark}[Cosimplicial vs.\ simplicial]
The dual category $\Delta^{\mathrm{op}}$ has the same objects but reverses
arrows. A \term{simplicial object} in a category $\mathcal{C}$ is a
functor $X : \Delta^{\mathrm{op}} \to \mathcal{C}$ — equivalently, a
sequence of objects $X_0, X_1, \ldots$ together with maps $d_i =
X(\delta^i) : X_n \to X_{n-1}$ (\emph{face} operators on the simplicial
side) and $s_i = X(\sigma^i) : X_n \to X_{n+1}$ (\emph{degeneracy}
operators), satisfying the \term{simplicial identities} — the relations
of Theorem~\ref{thm:cosimplicial-identities} with arrows reversed and
indices renamed. Chapter~38 develops this.
\end{remark}


% ═══════════════════════════════════════════════════════════════════════════
\section{The Universal Property of $\Delta$}
\label{sec:universal}
% ═══════════════════════════════════════════════════════════════════════════

\subsection{Informally}

We have seen that $\Delta$ has a clean generators-and-relations
presentation. There is also a much deeper structural reason for $\Delta$
to be ``the right'' category for indexing simplicial data: $\Delta$ is
the \emph{free} object of a particular categorical type. More precisely:
its augmented variant $\Delta_+$ is the free strict monoidal category
on a single monoid object. This means: \emph{specifying a monoid in
\emph{any} strict monoidal category is the same as specifying a strict
monoidal functor out of $\Delta_+$}.

The consequence: anywhere a monoid appears (and monoids are everywhere
in mathematics), the data is automatically organized by $\Delta$. This
is the structural origin of the omnipresence of simplicial objects.

\begin{readernote}
The universal property in this section is Joyal's theorem
(Theorem~\ref{thm:joyal-universal}). It is one of the most important
theorems of combinatorial category theory, but it requires the
augmented variant $\Delta_+$ — the version of $\Delta$ with the empty
ordered set $[-1] = \emptyset$ added as an initial object. The empty
ordinal provides the \emph{unit} for ordinal sum. Without it, $\Delta$
would not be strict-monoidal.
\end{readernote}

\subsection{Formalizing}

\begin{definition}[Augmented simplex category $\Delta_+$]
\label{def:augmented-delta}
The \term{augmented simplex category} $\Delta_+$ has the same data as
$\Delta$ together with an additional object
\begin{equation*}
  [-1] \;=\; \emptyset,
\end{equation*}
the empty ordered set. There is exactly one morphism $[-1] \to [n]$ for
each $[n]$ (the empty function), and no morphisms $[n] \to [-1]$ for
$n \geq 0$ (since $[n]$ is non-empty and $[-1]$ is empty).
\end{definition}

\begin{nameorigin}
\term{Augmented} from Latin \emph{augere} (``to increase, to enlarge'').
$\Delta_+$ is the ``augmented'' simplex category because we have
\emph{augmented} $\Delta$ by adding an initial object below it. The
notation $\Delta_+$ is older than current conventions for general
augmentations; some references write $\Delta_a$ or $\Delta_{\leq}$.
\end{nameorigin}

\begin{howtosay}
$\Delta_+$ \sayit{``$\Delta$-plus'' or ``augmented Delta''}; $[-1]$
\sayit{``the empty ordinal'' or ``negative-one bracket''}. The phrase
``augmented simplex category'' is standard in modern higher category
theory; older topology texts sometimes call it the ``simplicial category
with augmentation.''
\end{howtosay}

\begin{definition}[Ordinal sum]
\label{def:ordinal-sum}
The \term{ordinal sum} on $\Delta_+$ is the functor
$\oplus : \Delta_+ \times \Delta_+ \to \Delta_+$ given on objects by
\begin{equation*}
  [m] \oplus [n] \;=\; [m + n + 1],
\end{equation*}
with the convention that $[-1] \oplus [n] = [n] = [n] \oplus [-1]$. On
morphisms $f : [m] \to [m']$ and $g : [n] \to [n']$, the morphism
$f \oplus g : [m + n + 1] \to [m' + n' + 1]$ acts as $f$ on the first
$m + 1$ elements (sending $i \mapsto f(i)$) and as a shifted $g$ on the
remaining $n + 1$ elements (sending $m + 1 + j \mapsto m' + 1 + g(j)$).
\end{definition}

\begin{nameorigin}
The \term{ordinal sum} is the standard set-theoretic operation on
well-ordered sets: place $[m]$ first, then $[n]$ after — with every
element of $[m]$ less than every element of $[n]$. The result is the
ordered set of cardinality $(m+1) + (n+1) = m + n + 2$, which is
$[m + n + 1]$ in our indexing. Note that $\oplus$ is \emph{not}
symmetric: $[m] \oplus [n]$ has $[m]$'s elements coming first, but as
ordered sets is canonically isomorphic to $[n] \oplus [m]$.
\end{nameorigin}

\begin{theorem}[Strict monoidal structure on $\Delta_+$]
\label{thm:delta-plus-monoidal}
The functor $\oplus$ makes $\Delta_+$ a \emph{strict monoidal category}
with unit $[-1]$. Associativity is strict:
\begin{equation*}
  ([m] \oplus [n]) \oplus [p] \;=\; [m + n + p + 2] \;=\; [m] \oplus ([n] \oplus [p]).
\end{equation*}
\end{theorem}

\begin{remark}[A monoid object in $\Delta_+$]
The object $[0]$ carries a canonical monoid structure in the strict
monoidal category $(\Delta_+, \oplus, [-1])$. The structure maps are:
\begin{itemize}
  \item Multiplication: $\mu = \sigma^0 : [0] \oplus [0] = [1] \to [0]$
        (the unique monotone surjection).
  \item Unit: $\eta : [-1] \to [0]$ (the unique empty-function).
\end{itemize}
The unit and associativity laws hold strictly by the uniqueness of
monotone maps with prescribed targets in $\Delta_+$.
\end{remark}

\subsection{Reorganizing: Joyal's theorem}

The remarkable fact is that this single monoid object \emph{generates}
$\Delta_+$ in a precise sense.

\begin{nameorigin}
\term{Joyal's universal property of $\Delta_+$} is attributed to
\emph{André Joyal} (born 1943, Montréal), a foundational figure in
modern category theory. Joyal's contributions include species
(combinatorial structures), the Joyal model structure on simplicial sets
(making $\mathbf{sSet}$ a model for $\infty$-categories — Chapter~44),
and many results in higher categorical algebra. The universal property
theorem here is one of the cleanest applications of the
``free-monoid-on-an-object'' viewpoint in classical category theory.
\end{nameorigin}

\begin{theorem}[Joyal's universal property of $\Delta_+$]
\label{thm:joyal-universal}
The pair $(\Delta_+, [0])$ is the \term{free strict monoidal category on
a monoid object}: for any strict monoidal category $(\mathcal{V},
\otimes, I)$ equipped with a monoid object $(M, \mu : M \otimes M \to M,
\eta : I \to M)$, there is a unique strict monoidal functor
\begin{equation*}
  F : \Delta_+ \to \mathcal{V} \quad \text{with} \quad F([0]) = M, \;
  F(\sigma^0_{[1]}) = \mu, \; F(\eta_{[0]}) = \eta.
\end{equation*}
\end{theorem}

\begin{proof}[Proof sketch]
On objects, $F$ is forced: $F([n]) = M^{\otimes (n+1)}$ and $F([-1]) = I$.
On morphisms, the strict monoidal extension is unique because every
morphism in $\Delta_+$ factors as a composite of $\oplus$-tensorings of
the elementary morphisms $\sigma^0 : [1] \to [0]$ (the multiplication on
$[0]$), $[-1] \to [0]$ (the unit), and their identities — and these are
forced to map to $\mu$, $\eta$, and identities. The cosimplicial
identities translate exactly into the monoid axioms, so $F$ is
well-defined.
\qed
\end{proof}

\begin{readernote}
Joyal's theorem is the \emph{universal explanation} for why $\Delta$
appears in so many contexts:
\begin{itemize}
  \item Wherever a monoid acts (and monoids are everywhere — groups,
        rings, monads, $A_\infty$-algebras, etc.), there is automatically
        a strict monoidal functor out of $\Delta_+$.
  \item Restricting that functor to $\Delta \subset \Delta_+$ produces a
        simplicial object — the canonical \emph{bar construction}
        (Corollary~\ref{cor:bar-construction}).
  \item This is a single sentence that explains the bar resolution in
        homological algebra, the nerve of a category, the simplicial
        resolution of a group, the bar-cobar duality of operads (Chapter~57),
        and the $\infty$-monadic bar (Chapter~51).
\end{itemize}
Internalize this and you will see $\Delta$ wherever you look.
\end{readernote}

\subsection{Formally: Consequences}

\begin{corollary}[Bar construction]
\label{cor:bar-construction}
For any monoid $M$ in a strict monoidal category $\mathcal{V}$, the
corresponding strict monoidal functor $F : \Delta_+ \to \mathcal{V}$
restricted to $\Delta \subset \Delta_+$ defines a simplicial object
\begin{equation*}
  B_\bullet M : \Delta^{\mathrm{op}} \to \mathcal{V}, \quad
  B_n M \;=\; M^{\otimes (n+1)},
\end{equation*}
known as the \term{bar construction} on $M$. Its face maps come from
multiplication; its degeneracy maps come from the unit.
\end{corollary}

\begin{nameorigin}
The \term{bar construction} gets its name from the vertical-bar notation
used by Eilenberg and Mac Lane (1953) for its elements: they wrote
$[a_1 | a_2 | \cdots | a_n]$ for an element of $M^{\otimes n}$, with
vertical bars separating the tensor factors. The construction itself
predates the notation by years, but the bars stuck.
\end{nameorigin}

\begin{remark}[Bar construction in Volume~IV's arc]
The bar construction appears repeatedly in Volume~IV. In Chapter~42 it
provides resolutions for derived functors; in Chapter~51 it computes the
universal $\infty$-monad; in Chapter~56 it relates $A_\infty$- to
$\mathbb{E}_1$-algebras; in Chapter~57 it is the keystone of $A_\infty$/
$L_\infty$ Koszul duality. Each appearance is a special case of the
construction here, applied in a different monoidal category.
\end{remark}

\begin{corollary}[The epi-mono factorization revisited]
\label{cor:epi-mono}
Theorem~\ref{thm:delta-generators}'s unique factorization is the
\emph{orthogonal factorization system} on $\Delta$ generated by
$\mathcal{E} = $ surjective monotone maps (composites of degeneracies)
and $\mathcal{M} = $ injective monotone maps (composites of faces).
Every monotone map factors uniquely (up to canonical isomorphism) as
$\mathcal{E}$ followed by $\mathcal{M}$.
\end{corollary}


% ═══════════════════════════════════════════════════════════════════════════
\section{Geometric Realization (Preview)}
\label{sec:geometric-realization-preview}
% ═══════════════════════════════════════════════════════════════════════════

\subsection{Informally}

We have been carrying a geometric mental picture all along — $[n]$ as
the vertex set of an $n$-simplex. The picture is more than a metaphor:
each object $[n] \in \Delta$ has a literal topological avatar, and the
morphisms of $\Delta$ encode literal continuous maps between simplices.

\begin{nameorigin}
\term{Geometric realization} is so named because it ``realizes'' the
combinatorial simplex $[n]$ as a literal geometric object in $\mathbb{R}^{n+1}$.
The verb \emph{realize} in mathematics typically means ``construct a
concrete instance of''; the noun \emph{realization} (or \emph{realisation},
British) names that instance. The terminology dates to early algebraic
topology and was standardized by Milnor's \emph{Geometric realization of
a semi-simplicial complex} (1957).
\end{nameorigin}

\subsection{Formally}

\begin{definition}[Standard topological $n$-simplex]
\label{def:topological-simplex}
The \term{standard topological $n$-simplex} is
\begin{equation*}
  \lvert\Delta^n\rvert \;=\; \Bigl\{ (t_0, \ldots, t_n) \in \mathbb{R}^{n+1}
  \;\Big|\; t_i \geq 0, \ \sum_i t_i = 1 \Bigr\}.
\end{equation*}
The $i$-th vertex is the standard basis vector $e_i \in \mathbb{R}^{n+1}$;
the simplex itself is the convex hull of $e_0, e_1, \ldots, e_n$.
\end{definition}

\begin{howtosay}
$\lvert\Delta^n\rvert$ \sayit{``the geometric $n$-simplex'' or ``the
realization of $\Delta^n$''}. The vertical bars $\lvert{-}\rvert$ are the
standard notation for realization throughout algebraic topology; we use
them consistently for the realization functor of Chapter~39.
\end{howtosay}

\begin{remark}[The cosimplicial space $\lvert\Delta^\bullet\rvert$]
The assignment $[n] \mapsto \lvert\Delta^n\rvert$ defines a functor
$\lvert\Delta^\bullet\rvert : \Delta \to \mathbf{Top}$, the
\term{cosimplicial space of standard simplices}. Face maps $\delta^i$
correspond to inclusions of $(n-1)$-faces into the $n$-simplex
(literally: send $e_j \mapsto e_{\delta^i(j)}$ and extend affinely).
Degeneracy maps $\sigma^i$ correspond to collapses of $n$-simplices onto
their lower-dimensional faces (send $e_j \mapsto e_{\sigma^i(j)}$ and
extend affinely; this collapses an edge of the $n$-simplex onto a
vertex of the $(n-1)$-simplex).
\end{remark}

\begin{theorem}[$\lvert\Delta^\bullet\rvert$ is fully faithful]
\label{thm:delta-faithful}
The functor $\lvert\Delta^\bullet\rvert : \Delta \to \mathbf{Top}$ is
faithful, and its image consists of the standard simplices and the
affine simplicial maps between them.
\end{theorem}

The functor $\lvert\Delta^\bullet\rvert$ is the engine of \term{geometric
realization}, which Chapter~39 extends along the Yoneda embedding
$\Delta \hookrightarrow \mathbf{sSet}$ to a left adjoint
$\lvert{-}\rvert : \mathbf{sSet} \to \mathbf{Top}$. The full development
is the content of Chapter~39; here we only note the existence as
motivation for the indexing combinatorics.

\begin{readernote}
The functor $\lvert\Delta^\bullet\rvert$ is not full — there are continuous
maps $\lvert\Delta^m\rvert \to \lvert\Delta^n\rvert$ that are \emph{not}
affine (e.g., any non-affine continuous map). Its faithfulness (in the
sense of injection on morphisms) is what matters for combinatorial
applications: distinct monotone maps yield distinct affine simplicial
maps.
\end{readernote}


% ═══════════════════════════════════════════════════════════════════════════
\section{Exercises}
\label{sec:simplex-exercises}
% ═══════════════════════════════════════════════════════════════════════════

\begin{exercisesbox}
\begin{qa}
\qaitem{What are the objects of $\Delta$?}{Finite non-empty totally ordered sets, presented as $[n] = \{0 < 1 < \cdots < n\}$ for $n \geq 0$. Up to canonical isomorphism, one per positive cardinality.}
\qaitem{What are the morphisms?}{Monotone (order-preserving) functions: $f : [m] \to [n]$ with $i \leq j \Rightarrow f(i) \leq f(j)$.}
\qaitem{Why does $[n]$ have $n+1$ elements, not $n$?}{Geometric convention: $[n]$ is the combinatorial $n$-simplex, which is an $n$-dimensional shape with $n+1$ vertices.}
\qaitem{Why ordered sets and not just sets?}{The ordering gives simplices orientation: a canonical labeling of vertices, so face maps and degeneracies can be defined unambiguously. $\Delta$ has no non-identity automorphisms — every $[n]$ is rigid.}
\qaitem{Where does the name ``simplex'' come from?}{Latin \emph{simplex}, ``simple'' — an $n$-simplex is the simplest convex shape of dimension $n$.}
\qaitem{Is $[0]$ initial, terminal, both, or neither in $\Delta$?}{Terminal: from any $[n]$ there is exactly one monotone map to $[0]$ (the constant). Not initial: e.g., two distinct monotone maps $[0] \to [1]$ exist.}
\qaitem{How many monotone maps $[m] \to [n]$ are there?}{$\binom{m+n+1}{m+1}$ — equivalently, the number of multisets of size $m+1$ from an $(n+1)$-element set.}
\qaitem{What is the face map $\delta^i : [n-1] \to [n]$?}{The unique injective monotone map whose image omits exactly $i \in [n]$: $\delta^i(j) = j$ if $j < i$, else $j+1$.}
\qaitem{What is the degeneracy map $\sigma^i : [n+1] \to [n]$?}{The unique surjective monotone map whose fibre over $i$ has two elements: $\sigma^i(j) = j$ if $j \leq i$, else $j-1$.}
\qaitem{Why are the maps named ``face'' and ``degeneracy''?}{``Face'' is geometric: $\delta^i$ embeds the $(n-1)$-simplex as the face opposite vertex $i$. ``Degeneracy'' from Latin \emph{degenerare}: $\sigma^i$ degenerates an $(n+1)$-simplex by collapsing two adjacent vertices.}
\qaitem{State one cosimplicial identity.}{$\sigma^j \delta^i = \mathrm{id}$ if $i \in \{j, j+1\}$ — doubling a point and then deleting one copy is the identity.}
\qaitem{What does $\delta^j \delta^i = \delta^i \delta^{j-1}$ for $i < j$ say geometrically?}{Two non-adjacent face inclusions of an $(n-2)$-face into an $n$-face commute, with the second-applied face index shifted to account for the absence of the first.}
\qaitem{State the generation theorem.}{Every morphism in $\Delta$ factors uniquely as $\delta^{i_1} \cdots \delta^{i_p} \, \sigma^{j_1} \cdots \sigma^{j_q}$ with $i_1 > \cdots > i_p$ and $j_1 < \cdots < j_q$.}
\qaitem{What is the epi-mono factorization in $\Delta$?}{Every monotone map factors as a surjection (degeneracies) followed by an injection (faces). The factorization is unique.}
\qaitem{What is the augmented simplex category $\Delta_+$?}{$\Delta$ with the empty ordered set $[-1] = \emptyset$ added as an object; the empty function $[-1] \to [n]$ is the unique morphism from $[-1]$.}
\qaitem{Why is $[-1]$ included in $\Delta_+$?}{So that ordinal sum has a unit, making $\Delta_+$ a strict monoidal category. In $\Delta$ alone, no object plays that role.}
\qaitem{What does ``augmented'' mean here?}{From Latin \emph{augere} (``to increase''); $\Delta_+$ is $\Delta$ augmented by an initial object.}
\qaitem{Define ordinal sum $[m] \oplus [n]$.}{$[m] \oplus [n] = [m + n + 1]$, with $[-1]$ as the unit. The first $m+1$ elements come from $[m]$; the last $n+1$ come from $[n]$ shifted.}
\qaitem{What monoid object lives in $\Delta_+$?}{$[0]$, with multiplication $[0] \oplus [0] = [1] \to [0]$ (the unique surjection $\sigma^0$) and unit $[-1] \to [0]$ (the empty function).}
\qaitem{State Joyal's universal property.}{$(\Delta_+, [0])$ is the free strict monoidal category on a monoid object: strict monoidal functors out of $\Delta_+$ correspond bijectively to monoids in the target monoidal category.}
\qaitem{Who is André Joyal?}{Born 1943, Montréal. Foundational figure in modern category theory; introduced species (combinatorial structures), the Joyal model structure on $\mathbf{sSet}$ (Chapter~44), and many results in higher categorical algebra. Theorem~\ref{thm:joyal-universal} is one of his cleanest results.}
\qaitem{What is the bar construction $B_\bullet M$?}{The simplicial object $\Delta^{\mathrm{op}} \to \mathcal{V}$ with $B_n M = M^{\otimes(n+1)}$ obtained by restricting Joyal's monoid-classifying functor to $\Delta \subset \Delta_+$.}
\qaitem{Where does the ``bar'' in ``bar construction'' come from?}{The vertical bars in the original notation $[a_1 | a_2 | \cdots | a_n]$ for elements; Eilenberg–MacLane introduced the notation in 1953. The bars separated tensor factors.}
\qaitem{Where does the bar construction appear in Volume IV?}{In resolutions for derived functors (Ch.~42); in the universal $\infty$-monad (Ch.~51); in relating $A_\infty$- to $\mathbb{E}_1$-algebras (Ch.~56); in Koszul duality (Ch.~57).}
\qaitem{What is the standard topological $n$-simplex?}{$\lvert\Delta^n\rvert = \{(t_0, \ldots, t_n) \in \mathbb{R}^{n+1} : t_i \geq 0, \sum t_i = 1\}$: the convex hull of the standard basis vectors.}
\qaitem{What does the functor $\lvert\Delta^\bullet\rvert : \Delta \to \mathbf{Top}$ do on morphisms?}{It extends a monotone $f : [m] \to [n]$ affinely: the $i$-th vertex of $\lvert\Delta^m\rvert$ goes to the $f(i)$-th vertex of $\lvert\Delta^n\rvert$, interpolated linearly.}
\qaitem{Why does $\Delta$ appear ``everywhere a monoid acts''?}{Because Joyal's theorem identifies $\Delta_+$ as the universal monoid-classifying monoidal category: every monoid in any strict monoidal category factors through $\Delta_+$.}
\qaitem{Is $\Delta$ small? Locally finite?}{Small (countably many objects) and locally finite (finite hom-sets, by the binomial count).}
\qaitem{What is the relationship between $\Delta$ and $\Delta^{\mathrm{op}}$?}{Opposite categories: $\Delta^{\mathrm{op}}$ has the same objects but reverses morphisms. A presheaf $\Delta^{\mathrm{op}} \to \mathbf{Set}$ is a simplicial set (Chapter~38).}
\qaitem{Why are the operators on a simplicial set written $d_i, s_i$ (subscripts) instead of $\delta^i, \sigma^i$ (superscripts)?}{To signal the direction: $\delta^i, \sigma^i$ live in $\Delta$ (covariant); $d_i = X(\delta^i)$ and $s_i = X(\sigma^i)$ are induced by a contravariant functor $X : \Delta^{\mathrm{op}} \to \mathbf{Set}$, hence go ``the other way.''}
\end{qa}
\end{exercisesbox}


% ═══════════════════════════════════════════════════════════════════════════
\section*{Chapter Summary}
\addcontentsline{toc}{section}{Chapter Summary}
% ═══════════════════════════════════════════════════════════════════════════

\begin{summarybox}
\textbf{The category $\Delta$.}\quad Objects are the finite non-empty
totally ordered sets $[n] = \{0 < 1 < \cdots < n\}$. Morphisms are
monotone (order-preserving) functions. $[0]$ is terminal; the hom-sets
are finite ($\binom{m+n+1}{m+1}$ monotone maps from $[m]$ to $[n]$).

\textbf{Faces and degeneracies.}\quad The face map $\delta^i : [n-1]
\to [n]$ is the injection skipping $i$; the degeneracy $\sigma^i :
[n+1] \to [n]$ is the surjection doubling $i$. Together they generate
every morphism: every $f : [m] \to [n]$ factors uniquely as $\delta^{i_1}
\cdots \delta^{i_p} \sigma^{j_1} \cdots \sigma^{j_q}$ in normal form.

\textbf{Cosimplicial identities.}\quad Five relations governing how
$\delta$'s and $\sigma$'s commute past each other or annihilate. They
encode exactly the relations among monotone maps.

\textbf{The augmented simplex category $\Delta_+$.}\quad Add the empty
ordered set $[-1]$ to make ordinal sum $[m] \oplus [n] = [m + n + 1]$
into a strict monoidal product with unit $[-1]$.

\textbf{Joyal's theorem.}\quad $\Delta_+$ is the free strict monoidal
category on a monoid object: strict monoidal functors $\Delta_+ \to
\mathcal{V}$ correspond bijectively to monoid objects in $\mathcal{V}$.
The bar construction $B_\bullet M$ is the resulting simplicial object.

\textbf{Geometric realization (preview).}\quad The functor $\Delta \to
\mathbf{Top}$ sending $[n]$ to the standard $n$-simplex
$\lvert\Delta^n\rvert$ is faithful and extends along $\Delta \hookrightarrow
\mathbf{sSet}$ to the geometric realization of Chapter~39.
\end{summarybox}


% ═══════════════════════════════════════════════════════════════════════════
\section*{Formal Restatement}
\addcontentsline{toc}{section}{Formal Restatement}
% ═══════════════════════════════════════════════════════════════════════════

\begin{formalrestatement}
\textbf{Simplex category.}\quad $\Delta = \{[n] : n \geq 0\}$, objects
$[n] = \{0 < 1 < \cdots < n\}$; morphisms are monotone functions.
$\Delta$ is small, locally finite, with $[0]$ terminal.

\medskip
\textbf{Face map.}\quad $\delta^i : [n-1] \to [n]$, $\delta^i(j) = j$ if
$j < i$, else $j+1$. Injective; image omits $i$.

\medskip
\textbf{Degeneracy map.}\quad $\sigma^i : [n+1] \to [n]$, $\sigma^i(j) = j$
if $j \leq i$, else $j-1$. Surjective; fibre over $i$ has two elements.

\medskip
\textbf{Cosimplicial identities.}\quad
\begin{align*}
  \delta^j \delta^i &= \delta^i \delta^{j-1}, & i &< j; \\
  \sigma^j \sigma^i &= \sigma^i \sigma^{j+1}, & i &\leq j; \\
  \sigma^j \delta^i &= \delta^i \sigma^{j-1}, & i &< j; \\
  \sigma^j \delta^i &= \mathrm{id},           & i &\in \{j, j+1\}; \\
  \sigma^j \delta^i &= \delta^{i-1} \sigma^j, & i &> j+1.
\end{align*}

\medskip
\textbf{Generators.}\quad Every $f \in \Delta([m], [n])$ has a unique
normal form $f = \delta^{i_1} \cdots \delta^{i_p} \sigma^{j_1} \cdots
\sigma^{j_q}$ with $i_1 > \cdots > i_p$ and $j_1 < \cdots < j_q$.
Equivalently: $\Delta$ admits a unique orthogonal factorization system
$(\text{surjections}, \text{injections})$.

\medskip
\textbf{Augmented simplex category.}\quad $\Delta_+ = \Delta \cup \{[-1]\}$
with $[-1] = \emptyset$ and unique empty maps $[-1] \to [n]$.

\medskip
\textbf{Ordinal sum.}\quad $\oplus : \Delta_+ \times \Delta_+ \to \Delta_+$,
$[m] \oplus [n] = [m+n+1]$, with $[-1]$ as strict unit. Strictly
associative. $(\Delta_+, \oplus, [-1])$ is a strict monoidal category.

\medskip
\textbf{Joyal's universal property.}\quad For any strict monoidal
$(\mathcal{V}, \otimes, I)$ and monoid object $(M, \mu, \eta)$ in
$\mathcal{V}$, there is a unique strict monoidal functor $F : \Delta_+
\to \mathcal{V}$ with $F([0]) = M$, $F(\sigma^0_{[1]}) = \mu$,
$F(\eta_{[0]}) = \eta$.

\medskip
\textbf{Bar construction.}\quad For $M \in \mathrm{Mon}(\mathcal{V})$,
the restriction of Joyal's $F$ to $\Delta^{\mathrm{op}}$ gives the
simplicial object $B_\bullet M$, $B_n M = M^{\otimes(n+1)}$, with
face/degeneracy operators induced by $\mu$/$\eta$.

\medskip
\textbf{Geometric realization (preview).}\quad $\lvert\Delta^n\rvert =
\{(t_0, \ldots, t_n) \in \mathbb{R}^{n+1} : t_i \geq 0, \sum t_i = 1\}$;
$\lvert\Delta^\bullet\rvert : \Delta \to \mathbf{Top}$ faithful. Extends
along Yoneda to $\lvert{-}\rvert : \mathbf{sSet} \to \mathbf{Top}$
(Chapter~39).
\end{formalrestatement}
